\section{Separated interpolation sets in the actual spectral interval}
\label{cg-section}

We now supply the separated interpolation sets left open in
Section~\ref{ca-section}.  The construction below is the one
used in the source: Chebyshev nodes in disjoint intervals,
together with a separate set in a shorter interval near the
left endpoint~\cite{alpo-uniform-source}.  Keeping the exact
Chebyshev coefficient sum improves the constant $28$ to $24$
at every degree needed in the argument.

\subsection{The spectral interval}

Use the parameters and polynomials $P,Q$ of
Section~\ref{as-section}.  Fix
$1\leq l\leq k-2$ with $\epsilon_l\ne\epsilon'_l$ and put
\[
 \alpha=\frac52-v_l,\qquad
 \beta=\frac{\Delta}{2}-v_l,\qquad
 I=\beta-\alpha=\frac{\Delta}{2}-\frac52.
\]
Let $r\geq0$ be an integer, $N=20(r+1)$ and $R=2r$.
For the construction assume
\begin{equation}\label{cg-length}
                        I\geq6N^3+6N.
\end{equation}
In particular $\Delta>5$.

\begin{lemma}\label{cg-spectral}
For every $Y\in[\alpha,\beta]$, $P(Y)\ne0$ and $Q(Y)>0$.
Also $-Y$ lies between the two differing spectral entries
with a margin of two:
\begin{equation}\label{cg-OnCut}
 \min(u_{l+1}(\epsilon),u_{l+1}(\epsilon'))+2
       \leq -Y
       \leq\max(u_{l+1}(\epsilon),u_{l+1}(\epsilon'))-2.
\end{equation}
\end{lemma}

\begin{proof}
Write $\tau=v_l+Y$, so $5/2\leq\tau\leq\Delta/2$.
The factors $Y+v_i=\tau+4\Delta(l-i)$ are positive
for $i\leq l$ and negative for $i>l$.
Furthermore
\[
 Y=\tau-v_l\leq\Delta/2-v_*-4\Delta(k-1-l)<0,
\]
so every $Y-b_j$ is negative.  This proves $P(Y)\ne0$.
The factor $Y+1/2+a$ is positive by
$a-v_1\geq4k\Delta$.
For $1\leq i\leq k-2$, the two factors of $Q$
associated with $u_{i+1}$ are
\[
              \tau+4\Delta(l-i)-\epsilon_i\Delta,
       \qquad \tau+4\Delta(l-i)-\epsilon'_i\Delta.
\]
They are both positive if $i<l$, and both negative if
$i>l$.  At $i=l$ they are $\tau$ and $\tau-\Delta$,
in some order, so their product is negative.
The two factors associated with $u_k$ are equal and
nonzero.  In the first case their common value is
$Y+v_*<0$; in the second it is $Y+1/2-b_*<0$.
Thus the product defining $Q$ has exactly one negative
paired factor.  Its initial factor $-TT'$ is negative,
so $Q(Y)>0$.

The two entries in \eqref{cg-OnCut} are
$v_l-1/2-\Delta$ and $v_l-1/2$.
The upper inequality is $\tau\geq5/2$.
The lower is $\tau\leq\Delta-3/2$, which follows from
$\tau\leq\Delta/2$ and $\Delta\geq3$.
\end{proof}

In particular $P(y)$ and $P(y+1)$ are nonzero whenever
$y,y+1\in[\alpha,\beta]$.  The same margin condition
holds at both parameters.

\subsection{The finite sets}

Put
\[
 q=\frac{I}{6N},\qquad \ell=q-1,\qquad
 A_a=\alpha+(2N+2a-2)q\quad(1\leq a\leq N).
\]
Condition~\eqref{cg-length} gives $q\geq N^2+1$ and
$\ell\geq N^2>0$.  Define
\[
                    J_a=[A_a,A_a+\ell].
\]
For $d\geq1$, let
\[
 c_{a,d,i}
   =A_a+\frac{\ell}{2}
                \left(1+\cos\frac{i\pi}{d}\right)
                       \quad(0\leq i\leq d),
\]
and set
\[
 \mathcal G_a=
       \bigcup_{1\leq d\leq\max(1,R)}
                    \{c_{a,d,0},\ldots,c_{a,d,d}\}.
\]
For positive $a$ outside $\{1,\ldots,N\}$ let
$\mathcal G_a=\varnothing$.
Independently choose arbitrary
$\delta_0,\ldots,\delta_{N-1}\in[0,1/2]$ and put
\[
 z_j=\alpha+jN^2+\delta_j\quad(0\leq j<N),\qquad
 \mathcal G_0=\{z_0,\ldots,z_{N-1}\}.
\]
All choices of the $\delta_j$ have the properties below.
The source's short-interval nodes, which additionally avoid
the finitely many zeros of $P(Y)P(Y+1)$ by construction,
are among these choices.  Here that avoidance also follows
from Lemma~\ref{cg-spectral}.

\begin{proposition}\label{cg-sets}
The set $\mathcal G_0$ has exactly $N$ elements.
Every $y$ in any of the displayed sets satisfies
$y,y+1\in[\alpha,\beta]$, and every $y\in\mathcal G_0$
also satisfies
\[
                 \frac52\leq v_l+y,\quad
                              v_l+y+1\leq N^3.
\]
For $e,e'\in\{0,1\}$, the following separations hold:
\begin{align}
 |(y+e)-(y'+e')|&\geq N^2+1
       &&(y\in\mathcal G_a,\ y'\in\mathcal G_{a'},\
                   1\leq a,a'\leq N,\ a\ne a'),\label{cg-long-sep}\\
 |(y+e)-(y'+e')|&\geq N^2+1
       &&(y\in\mathcal G_0,\ y'\in\mathcal G_a,\
                   1\leq a\leq N),\label{cg-cross-sep}\\
 |(y+e)-(y'+e')|&\geq N^2-\frac32\geq\frac{N^2}{8}
       &&(y,y'\in\mathcal G_0,\ y\ne y').\label{cg-short-sep}
\end{align}
\end{proposition}

\begin{proof}
Every long-interval point satisfies
\[
 \alpha+2Nq\leq y,\qquad
             y+1\leq\alpha+(4N-1)q<\alpha+6Nq=\beta.
\]
For a short-interval point,
\[
 \alpha\leq y,\qquad
 y+1\leq\alpha+(N-1)N^2+\frac32
                      \leq\alpha+N^3-\frac52.
\]
The last inequality uses $N^2\geq4$, and the last upper
bound is less than $\beta$ by \eqref{cg-length}.
It also proves the stated shorter-interval bound.

If $a<a'$, the least possible difference
$(y'+e')-(y+e)$ is at least
\[
             A_{a'}-(A_a+\ell+1)
                    =\bigl(2(a'-a)-1\bigr)q\geq q.
\]
This proves \eqref{cg-long-sep}.
The least difference between a long point with its shift
and a short point with its shift is at least
\[
          2Nq-N^3+\frac52
                 \geq N^3+2N+\frac52>N^2+1,
\]
proving \eqref{cg-cross-sep}.
If $j<j'$, then
$z_{j'}-z_j\geq N^2-1/2>0$.
Thus the $z_j$ are distinct, and allowing the two shifts
decreases their distance by at most one.
Finally $N^2-3/2\geq N^2/8$ for $N\geq20$.
\end{proof}

\begin{proposition}[Interpolation with the actual separation]
\label{cg-interpolation}
For every $1\leq a\leq N$ and every integer $0\leq d\leq R$,
there are $d+1$ distinct points $y_0,\ldots,y_d$ in
$\mathcal G_a$ such that their Lagrange leading-coefficient
weights satisfy
\begin{equation}\label{cg-weight}
 \sum_{i=0}^d\left|
       \prod_{h\ne i}(y_i-y_h)^{-1}\right|
                       \leq\left(\frac{24N}{I}\right)^d.
\end{equation}
\end{proposition}

\begin{proof}
For $d=0$, choose the endpoint $A_a+\ell$, which occurs
in the degree-one set.  The empty product is one, so
\eqref{cg-weight} holds.
For $d\geq1$ use $y_i=c_{a,d,i}$.
Lemma~\ref{ca-Chebyshev} gives the exact coefficient sum
$\frac12(4/\ell)^d$.  Since $q=I/(6N)$ and
$q\geq N^2+1$,
\[
 \frac4\ell=\frac{24N}{I}\frac{q}{q-1}
              \leq\frac{24N}{I}\left(1+\frac1{N^2}\right).
\]
For an integer $d\geq1$ and $x\geq0$ with $dx<1$,
\[
 (1+x)^d
     =\sum_{h=0}^d\binom dh x^h
     \leq\sum_{h=0}^{\infty}(dx)^h
     =\frac1{1-dx}.
\]
Here $d\leq R=2r<N/10$, so in particular
$d/N^2\leq1/2$ and the last quantity is at most two.
The factor $1/2$ in the exact coefficient sum absorbs
this loss.  This proves \eqref{cg-weight}.
The same proof works for any $d\leq N^2/2$, if the set
is constructed with nodes of that degree.
\end{proof}

\subsection{Extension of signed parameter sequences}

A signed parameter is a pair $(\varepsilon,t)$ with
$\varepsilon\in\{+,-\}$ and $t\in\mathbb R$.
Its underlying node is $t$ for sign $+$ and $t-1$
for sign $-$.
Call a finite sequence
$w=((\varepsilon_1,t_1),\ldots,(\varepsilon_h,t_h))$
admissible if $h\leq N$, its $i$th underlying node belongs to
$\mathcal G_{h+1-i}\cup\mathcal G_0$, and
\begin{equation}\label{cg-admissible}
       |t_i-t_j+\delta|>2
       \quad(i\ne j,\ \delta\in\mathbb Z,\ |\delta|\leq2).
\end{equation}
These are the actual per-place conditions used in the
source.  Every parameter belongs to $[\alpha,\beta]$
and hence satisfies \eqref{cg-OnCut}.

For a list $\boldsymbol y=(y_1,\ldots,y_K)$, write
\[
 w_+(\boldsymbol y)
       =((+,y_1),\ldots,(+,y_K),w),\qquad
 w_-(\boldsymbol y)
       =((-,y_1+1),\ldots,(-,y_K+1),w).
\]
The notation means concatenation with the original sequence.

\begin{theorem}\label{cg-extension}
Let $w$ be admissible of length $h$ and let $h+K\leq N$.
For every $0\leq d\leq R$ there are $K$ grids
$y_{a,0},\ldots,y_{a,d}$, $1\leq a\leq K$, satisfying
\eqref{cg-weight}, such that every selection
$\boldsymbol y=(y_{1,i_1},\ldots,y_{K,i_K})$
has the following properties:

\begin{enumerate}
\item Both $w_+(\boldsymbol y)$ and $w_-(\boldsymbol y)$
are admissible, have length $h+K$, and satisfy the
spectral margin \eqref{cg-OnCut} at every parameter.
\item The new nodes are pairwise at distance greater
than four.  For every old parameter $(\varepsilon,t)$,
every new node $y$ and every $e\in\{0,1\}$,
\[
                         |y+e-t|\geq N^2+1.
\]
\item The Lagrange product weights, normalized by their
sum of absolute values as in
Section~\ref{ca-section}, define a probability measure
on the resulting finite family of extended sequences.
\end{enumerate}
At degree $d=0$ one may instead choose $K$ distinct
nodes from $\mathcal G_0$ not already underlying old
parameters.  All conclusions still hold, with the
distance lower bound $N^2/8$ in place of $N^2+1$,
and every new node lies in the shorter interval above.
\end{theorem}

\begin{proof}
For $a=1,\ldots,K$ take the degree-$d$ grid from
$\mathcal G_{h+K+1-a}$ in
Proposition~\ref{cg-interpolation}.
These block indices are distinct and are all greater
than $h$.  Each old underlying node either belongs to
$\mathcal G_0$ or to a block with index at most $h$.
The separations \eqref{cg-long-sep} and
\eqref{cg-cross-sep} therefore apply to every new pair
and every new--old pair, for all shifts by zero or one.
In particular all those parameter distances exceed four.
Adding an integer of absolute value at most two leaves
a distance greater than two, proving
\eqref{cg-admissible} for the extended sequences.
The previously existing pairs already satisfy it.

For the block membership condition, the $a$th new
underlying node lies in
$\mathcal G_{h+K+1-a}$ by construction.
An old entry moves from position $i$ to $K+i$, and
\[
                  (h+K)+1-(K+i)=h+1-i.
\]
Thus its required block is unchanged.  Both signs
of extension use the same underlying nodes.  The
interval and spectral assertions follow from
Proposition~\ref{cg-sets} and Lemma~\ref{cg-spectral}.

Every Lagrange weight is nonzero, because the nodes
within its grid are distinct.  The product of the
normalized absolute weights sums to one by expansion
of a product of finite sums.  Assign the sum of these
weights to each sequence if several choices represent
the same sequence.  This gives the stated probability
measure without any counting factor.

For degree zero, at most $h$ elements of $\mathcal G_0$
are underlying nodes of old parameters.  Since
$|\mathcal G_0|=N\geq h+K$, at least $K$ unused
nodes remain.  Choose any $K$ of them.
Equation~\eqref{cg-short-sep} supplies all new--new
and relevant new--old separations within the short
set, and \eqref{cg-cross-sep} handles old long nodes.
Here $N^2/8\geq50>4$, so the same admissibility
argument applies.  Every grid has one node and weight
one.  For $K=0$ both constructions give the original
sequence and the probability measure of mass one on it.
\end{proof}

\begin{corollary}\label{cg-actual-constant}
The actual Fourier interpolation degree $d=2j$,
$0\leq j\leq r$, admits the construction above.
The denominators in \eqref{ca-q-distortion-eq} are
bounded below by $N^2/8$ for the short-interval
construction and by $N^2+1$ for the long one.
Consequently the algebraic barrier in
\eqref{ca-X-exact} may be replaced by
\[
 \frac{g\lambda}{2\pi}
        \left(\frac{24N}{I}\right)^{2j}
                             e^{10/49}A_{\rm pol}.
\]
This improvement retains the particular Hermitian-form
inequality and Fourier identity assumed in
Section~\ref{ca-section}.
\end{corollary}

\begin{proof}
For an old positive parameter, the denominator is
$y-t$; for an old negative parameter it is
$t-y-1$.  The latter is the difference between its
underlying node and $y$.  Both are covered by the
proved separation inequalities, using the appropriate
zero or one shift.  Apply
Lemma~\ref{ca-q-distortion} and
Proposition~\ref{ca-positive-block}.
\end{proof}

The interval condition \eqref{cg-length} also follows
eventually from the numerical range in
Section~\ref{as-section}.  Indeed, there $m=2r+2\leq k$,
so $N=10m\leq10k$.  If
$\Delta\geq(kf\log B)^{C_\Delta}$ with $C_\Delta>4$,
$f\geq2$ and $\log B\geq1$, then
$\Delta\geq k^{C_\Delta}$.
For all sufficiently large $k$, depending only on
$C_\Delta$, this gives
$\Delta/2-5/2\geq6000k^3+60k\geq6N^3+6N$.
Thus this construction introduces no new asymptotic
hypothesis in that range.

This closes the finite-set construction, all simultaneous
separation conditions, and the interpolation probabilities.
For a finite family of complex places, perform the same
construction at each place.  Its constants depend only on
$N,I$ and the degree, not on the translated interval
endpoint $\alpha$, the place, or the choices $\delta_j$.
The source bounds the lengths of the parameter sequences
by their maximum over the places.  Replacing the sequence
at one place by any extension of length at most $N$
therefore preserves the global length restriction.
It does not prove the radial Fourier identity, the
positivity and mixed inequality for the specified
differential forms, their wave isometry, or the analytic
domain and limiting assertions.  The separate mass,
source-dimension and formal-degree hypotheses remain
as stated earlier.


